\subsection{Оптические системы}

%%%%%%%%%%%%%%%%%%%%%%%%%%%%%%%%%%%%%%%%%%%%%%%%%%%%
%%%%%%%%%%%%%%%%%%%%%%%%%%%%%%%%%%%%%%%%%%%%%%%%%%%%
%%%%%%%%%%%%%%%%%%%%%%%%%%%%%%%%%%%%%%%%%%%%%%%%%%%%
 \z{На пути параллельного пучка лучей света ставят собирающую линзу, а за линзой экран. И линза и экран перпендикулярны оси пучка. На экране образуется светлая точка, если он находится на расстоянии $f = 30\,\text{см}$ от линзы. Вплотную к первой линзе приставляют вторую:  
	\lpu собирающую линзу с фокусным расстоянием $F = 20\,\text{см}$;   \lpuend
	\lpu рассеивающую линзу с фокусным расстоянием $F = 40\,\text{см}$.  \lpuend

\noindent
Где нужно поставить экран, чтобы на нем снова образовалась светлая точка?}
%
{\puan $l = 12\,\text{см}$; \quad $l = 120\,\text{см}$}

%%%%%%%%%%%%%%%%%%%%%%%%%%%%%%%%%%%%%%%%%%%%%%%%%%%%
 \z{Собирающая линза с оптической силой $D_1 = 2.45\,\text{дптр}$ сложена вплотную с рассеивающей линзой с фокусным расстоянием $F_2 = 0.5\,\text{м}$ так, что их оптические оси совпадают. Определить положение предмета, помещенного перед линзами, если его изображение отстоит от линз на расстоянии $f = 0.5\,\text{м}$. Какое получается изображение?}
%
{$d = 0.408\,\text{м}, \quad \text{мнимое}$}

%%%%%%%%%%%%%%%%%%%%%%%%%%%%%%%%%%%%%%%%%%%%%%%%%%%%
 \z{Две линзы с фокусным расстоянием $F_1 = 10\,\text{см}$ и $F_2 = 15\,\text{см}$ расположены на расстоянии $l = 30\,\text{см}$ друг от друга. Где следует расположить точечный источник света, чтобы идущие от него лучи после прохождения обеих линз образовали пучок, параллельный главной оптической оси?  
Задачу решить для случаев:  
\pu обе линзы собирающие;  
\lpu первая линза собирающая, вторая — рассеивающая.  \lpuend

\noindent
Решение проиллюстрировать построением. Можно ли получить при данном расстоянии между линзами параллельный пучок, если первая линза рассеивающая?
}
%
{\puan $d = 30\,\text{см}$; \quad $d = 12.8\,\text{см}$}

%%%%%%%%%%%%%%%%%%%%%%%%%%%%%%%%%%%%%%%%%%%%%%%%%%%%

 \z[{\begin{wrapfigure}[5]{r}{0.25\textwidth}%
		\vspace{-10pt}%
		\begin{flushright}%
 			\includegraphics[width=0.25\textwidth]{4.6.4}%
		\end{flushright}%
	\end{wrapfigure}%
	}]%
{Рассеивающая и собирающая линзы с одинаковыми по величине фокусными расстояниями $f = 10\,\text{см}$ расположены на расстоянии $f$ друг от друга, как показано на рисунке. Предмет находится на расстоянии $a = 20\,\text{см}$ от рассеивающей линзы. На каком расстоянии $b$ от собирающей линзы находится изображение предмета, показанное на рисунке штриховой линией?}
%
{$b = 2f + \frac{f^2}{a} = 25\,\text{см}$}

%%%%%%%%%%%%%%%%%%%%%%%%%%%%%%%%%%%%%%%%%%%%%%%%%%%%

 \z[{\begin{wrapfigure}[5]{r}{0.25\textwidth}%
		\vspace{-10pt}%
		\begin{flushright}%
 			\includegraphics[width=0.25\textwidth]{4.6.5}%
		\end{flushright}%
	\end{wrapfigure}%
	}]%
{Оптическая система состоит из двух собирающих линз, расположенных на расстоянии $l = 0.75\,\text{м}$ друг от друга. Наблюдателю, смотрящему справа, источник $S$, помещенный на главной оптической оси на расстоянии $d_1 = 0.25\,\text{м}$ от первой линзы, кажется удаленным на очень большое расстояние. Чему равна оптическая сила второй линзы, если оптическая сила первой линзы $D_2 = 6\,\text{дптр}$?}
%
{$D_2 = 4\,\text{дптр}$}

%%%%%%%%%%%%%%%%%%%%%%%%%%%%%%%%%%%%%%%%%%%%%%%%%%%%
 \z{Человек, страдающий близорукостью, рассматривает предмет, находящийся на расстоянии $d = 202\,\text{см}$ перед его глазами с использованием контактной линзы оптической силы $D = -5\,\text{дптр}$. При этом изображение предмета оказывается точно в плоскости сетчатки глаза. Определить, на какое расстояние $\delta$ сместится плоскость изображения, если человек снимет контактные линзы. Считать, что оптическая система глаза — это тонкая линза с фокусным расстоянием $f = 2\,\text{см}$, а контактная линза вплотную примыкает к ней.}
%
{$\delta = \frac{f d}{d - |D| f d - f} - \frac{f d}{d - f} \approx 2.2\,\text{мм} \quad \text{в сторону хрусталика}$}

%%%%%%%%%%%%%%%%%%%%%%%%%%%%%%%%%%%%%%%%%%%%%%%%%%%%
 \z{Оптическая система состоит из двух линз с фокусными расстояниями $F_1 = 20\,\text{см}$ и $F_2 = 40\,\text{см}$, преобразует пучок лучей, идущих от точечного источника, расположенного на главной оптической оси, на расстоянии $d_1 = 30\,\text{см}$ от первой линзы, в пучок лучей, параллельный главной оптической оси. На каком расстоянии друг от друга расположены линзы, если обе линзы собирающие?}
%
{$L = 20\,\text{см}$}

